\begin{abstract}
\noindent
\textbf{Motivation.} Phylogenetic network methods have accumulated faster than
the software infrastructure to support them. PhyloNet established a body of
network-analysis methods and the computational ecosystem around them, and much
of the field's methodological work now assumes that body of methods exists. New
methods, however, are typically built as standalone programs, each supplying its
own network representation, search machinery and input handling, and most are not
reachable from the Python environment where comparative and statistical analysis
is now usually done.

\noindent
\textbf{Results.} We present \phynetpy, a Python library for phylogenetic network
inference and analysis. Inference is expressed as two verbs, \code{infer} and
\code{score}, parameterised on three independent axes: the data observed, the
generative model, and the statistical criterion. A method is a cell in that
product rather than a command name, and dispatch runs through a registry that
doubles as a machine-readable statement of which combinations are defined,
which are unimplemented, and which are meaningless. Version 1.0.0 provides seven
inference methods, six of them reimplementations of published PhyloNet analyses;
a reusable rooted-DAG network type with ploidy and typed branch-length units;
nine network dissimilarity measures and a separate reticulation-specific
comparison module; coalescent, sequence and biallelic-marker simulation; and
MCMC convergence diagnostics with Tracer-compatible output. The library holds
1{,}130 tests, and a harness cross-validates its multispecies-network-coalescent
density and Felsenstein likelihood against PhyloNet's own implementations across
14 adversarial cases, agreeing to $10^{-5}$. On matched pseudo-likelihood
inference from identical gene trees, \phynetpy\ runs faster than PhyloNet by a
margin that grows with dataset size, but recovers the generating topology about
half as often across 30 trials; we report that discrepancy and the diagnostic
work it still needs rather than a net verdict.

\noindent
\textbf{Availability.} \code{pip install phynetpy}. Source, documentation and the
scripts reproducing every figure and table here are at
\url{https://github.com/NakhlehLab/PhyNetPy}, under the MIT license.
\end{abstract}
